\documentclass[10pt,a4paper]{amsart}
\usepackage[margin=19mm]{geometry}
\usepackage{amsmath,amssymb,mathtools}
\usepackage[scr=rsfs,cal=euler]{mathalfa}
\usepackage{xfrac}
\usepackage[unicode,hidelinks,bookmarks=false]{hyperref}
\newcommand{\ie}{i.e.\ }
\newcommand{\cf}{cf.\ }
\newcommand{\resp}{respectively\ }
\newcommand{\Subsection}{\subsection}
\providecommand{\nobreakdash}{\nobreak\hbox{-}}
\setlength{\emergencystretch}{2em}
\multlinegap0pt
\usepackage{tikz-cd}
\usepackage{mathrsfs}
\usepackage{changepage}
\usepackage{setspace}
\usepackage{fancyhdr}
\DeclareMathOperator{\Hom}{Hom}
\DeclareMathOperator{\Th}{Th}
\DeclareMathOperator{\End}{End}
\DeclareMathOperator{\Ima}{im}
\DeclareMathOperator{\id}{id}
\DeclareMathOperator{\Res}{Res}
\DeclareMathOperator{\Ind}{Ind}
\DeclareMathOperator{\tr}{tr}
\DeclareMathOperator{\Z}{\mathbb{Z}}
\DeclareMathOperator{\Q}{\mathbb{Q}}
\DeclareMathOperator{\R}{\mathbb{R}}
\DeclareMathOperator{\C}{\mathbb{C}}
\DeclareMathOperator{\N}{\mathbb{N}}
\DeclareMathOperator{\F}{\mathbb{F}}
\DeclareMathOperator{\irr}{irr}
\DeclareMathOperator{\TD}{tr.deg}
\DeclareMathOperator{\Aut}{Aut}
\DeclareMathOperator{\Tor}{Tor}
\DeclareMathOperator{\Ext}{Ext}
\DeclareMathOperator{\coker}{coker}
\DeclareMathOperator{\tp}{tp}
\DeclareMathOperator{\Spec}{Spec}
\DeclareMathOperator{\Proj}{Proj}
\DeclareMathOperator{\mfp}{\mathfrak{p}}
\DeclareMathOperator{\mfq}{\mathfrak{q}}
\DeclareMathOperator{\red}{red}
\DeclareMathOperator{\sHom}{\mathscr{H}\!om}
\DeclareMathOperator{\blank}{\underline{\enspace}}
\newcommand{\catname}[1]{{\normalfont\textbf{#1}}}
\DeclareMathOperator{\acl}{acl}
\DeclareMathOperator{\dcl}{dcl}
\def\Ind#1#2{#1\setbox0=\hbox{$#1x$}\kern\wd0\hbox to 0pt{\hss$#1\mid$\hss}
\lower.9\ht0\hbox to 0pt{\hss$#1\smile$\hss}\kern\wd0}
\def\ind{\mathop{\mathpalette\Ind{}}}
\def\notind#1#2{#1\setbox0=\hbox{$#1x$}\kern\wd0
\hbox to 0pt{\mathchardef\nn=12854\hss$#1\nn$\kern1.4\wd0\hss}
\hbox to 0pt{\hss$#1\mid$\hss}\lower.9\ht0 \hbox to 0pt{\hss$#1\smile$\hss}\kern\wd0}
\def\nind{\mathop{\mathpalette\notind{}}}
\newtheorem{theorem}{Theorem}[section]
\newtheorem{lemma}[theorem]{Lemma}
\newtheorem{prop}[theorem]{Proposition}
\newtheorem{remark}[theorem]{Remark}
\newtheorem{definition}[theorem]{Definition}
\newtheorem{exercise}[theorem]{Exercise}
\newtheorem*{exercise*}{Exercise}
\newtheorem{fact}[theorem]{Fact}
\begin{document}
\title[Some More Rigid Real Closed Fields]{Some More Rigid Real Closed Fields}
\author{Michael Lange ạte August 2026 e̱gin document 14pt 8pt; Michael Lange}
\date{}
\maketitle
\noindent\emph{Source and licence.} Some More Rigid Real Closed Fields. Version arXiv:2608.16766v1 (CC BY 4.0). This material is distributed under the Creative Commons Attribution 4.0 International licence, http://creativecommons.org/licenses/by/4.0/, as recorded in the arXiv OAI licence metadata for this exact version and shown on the abstract page https://arxiv.org/abs/2608.16766v1. Full rights notice: licenses/arxiv-2608.16766-CC-BY-4.0.txt.\par\smallskip
\noindent\textit{Editorial scope.} Counted unit: the original \texttt{theorem} environment \texttt{InfiniteThm} at source lines 150--152 together with its complete proof at lines 153--162; the delivered document keeps source lines 107--163 so that every referenced label and every citation resolves inside it. Native editable source is the paper's own concatenated TeX; the AMS wrapper is editorial formatting only. No publisher class, upstream script or unrelated artwork is redistributed.
\section{Constructing rigid fields}
A key fact for the main argument is that RCF (or any of its o-minimal expansions) has definable Skolem functions, which implies that for any subset $S$ of a model $N$, the definable closure of $S$ is an elementary submodel of $N$; this is clearly the smallest such submodel containing $S$. Given some elements $a_1,\dots,a_n\in R$ we use $k\langle a_1,\dots,a_n\rangle$ to denote this field, the real closure of $a_1,\dots,a_n$.
\\
\\
To produce a rigid non-Archimedean extension of finite transcendence degree over $k$, we employ exactly the same main argument as used in [1]. The extensions we construct will all have the form $k\langle a_1,\dots, a_n\rangle$ where $a_1>\alpha$ for all $\alpha\in k$ and all $a_i$ for $i\geq 2$ are in the convex hull of $k\langle a_1\rangle$ inside $R$. Given an automorphism $\sigma$ of such an extension $k\langle a_1,\dots, a_n\rangle$, we know that $\sigma(a_1),\dots,\sigma(a_n)\equiv a_1,\dots, a_n$. Also, since $k\langle a_1,\dots, a_n\rangle$ is the definable closure of $a_1,\dots,a_n$, there is some definable function $F_\sigma:R^n\to R^n$ such that $F_\sigma(a_1,\dots,a_n)=(\sigma(a_1),\dots,\sigma(a_n))$. If we wish to ensure that $k\langle a_1,\dots, a_n\rangle$ is rigid, it will suffice to ensure that every such definable map $F$ either fixes $(a_1,\dots,a_n)$ or else fails to preserve the type of this tuple. In the latter case, $F$ cannot be $F_\sigma$ for any automorphism $\sigma$. Hence, under our desired condition, for any automorphism $\sigma$, $F_\sigma$ must fix $(a_1,\dots,a_n)$ and so $\sigma$ must be the identity map.

To state the key lemma for this argument, we recall a definition from [1] (suitably adjusted for $n$ dimensions):
\begin{definition}
    An end-cell in $R^n$ is a definable set of the form 
    \[(\alpha,\infty)\times(g_2,h_2)\times\dots( g_n,h_n)\]
    where each $g_i,h_i$ is further assumed to be continuous on $(\alpha,\infty)\times\dots\times(g_{i-1},h_{i-1})$
\end{definition}

The main lemma we will need to complete this argument is the following generalization of Lemma 2.4 from [1]:

\begin{lemma}\label{MainLemma} Given a definable map $F:R^n\to R^n$ and an end cell $C$, there exists a possibly smaller end-cell $C'\subseteq C$ such that either $F\upharpoonright C'$ is the identity, or $F(C')\cap C'=\emptyset$. In fact, in the latter case we can ensure that the closure $\overline{C'}$ is disjoint from $F(C')$.
\end{lemma}

The additional claim about closure will not be needed until we reach the case of infinite transcendence degree. We will prove Lemma \ref{MainLemma} in the following section. Here we give the proof of the first main result using Lemma \ref{MainLemma}:

\begin{theorem}\label{FiniteThm}
    There exist non-Archimedean rigid real closed fields of all finite transcendence degrees $n\geq 2$ over $k$.
\end{theorem}
\begin{proof}
    Fix $n\geq 2$. Enumerate all definable maps $R^n\to R^n$ as $F_0,F_1,\dots$. Let $C_0:=(0,\infty)\times R^{n-1}$. Suppose $C_0\supseteq C_1\supseteq\dots\supseteq C_i$ have been defined such that $C_i$ is an end cell. Then let $C'_{i+1}$ be the end cell obtained by applying Lemma $\ref{MainLemma}$ to $C_i, F_i$ and take $C_{i+1}:=C'_{i+1}\cap \{i+1<x_1<\infty\}$. Now the intersection $\bigcap_{i<\omega} C_i$ may be viewed as a partial type $\tau(x_1,\dots,x_n)$ with the following properties: $\tau(x_1,\dots,x_n)\vdash x_1>\alpha$ for all $\alpha\in k$, and for any definable map $F:R^n\to R^n$, either $\tau(x_1,\dots,x_n)\vdash F(x_1,\dots,x_n)= (x_1,\dots,x_n)$ or else there is a formula $\varphi\in \tau$ such that $\varphi(x_1,\dots,x_n)\vdash \lnot\varphi(F(x_1,\dots,x_n))$. Namely, if $F=F_i$ under the chosen enumeration, then $\varphi$ may be taken as the definition of $C'_{i+1}$ in the above construction. We wish to take a realization $a_1,\dots,a_n$ of $\tau$ such that $a_1,\dots,a_n$ are algebraically independent. If this were not possible, then $\tau$ implies some algebraic relation among $a_1,\dots,a_n$ (in fact, one of the $a_i$ would be definable from the others). By compactness, some $C_i$ in the above construction implies this relation among $a_1,\dots,a_n$. This is not the case, however, as each $C_i$ has full dimension in $R^n$.
    
    Now consider the real closed field $K:=k\langle a_1,\dots,a_n\rangle$. This field is non-Archimedean as $a_1>\alpha$ for every $\alpha\in k$. Suppose $\sigma:K\to K$ is an automorphism. Then necessarily $\sigma(a_1),\dots,\sigma(a_n)\equiv a_1,\dots,a_n$. In particular, $\sigma(a_1),\dots,\sigma(a_n)\models \tau(x_1,\dots,x_n)$. Moreover, $\sigma(a_1),\dots,\sigma(a_n)\in k\langle a_1,\dots,a_n\rangle$ is exactly to say $(\sigma(a_1,\dots,\sigma(a_n))\in\text{dcl}(a_1,\dots,a_n)$, so there is a definable map $F:R^n\to R^n$ with $F(a_1,\dots,a_n)= (\sigma(a_1),\dots\sigma(a_n))$. Since $F(a_1,\dots,a_n)\models \tau(x_1,\dots,x_n)$, the construction of $\tau$ requires that $\tau(x_1,\dots,x_n)\vdash F(x_1,\dots,x_n)=(x_1,\dots,x_n)$. Then $\sigma$ is the identity on each $a_i$ and $K=\dcl(a_1,\dots,a_n)$ so $\sigma$ is the identity on $K$.
    
    
    
\end{proof}

As remarked in [1], one cannot directly employ the above argument to obtain a rigid non-Archimedean extension of degree $\omega$ by trying to diagonalize against all definable maps $R^\omega\to R^\omega$, as there are uncountably many such. This means that attempting to continue the above proof as a transfinite induction fails at the first limit step $\omega$: we can only possibly define $C_\omega$ to be some sub-cell of the intersection of all $C_i$ for $i<\omega$, but this intersection is not necessarily a definable set and so we have no guarantee of finding a sub-cell. It is essential that the collection we diagonalize against be countable. It will turn out that we can diagonalize against the countable collection of definable maps $R^n\to R^m$ for $n\geq m$.

Suppose we have an extension $K:=k\langle (a_i)_{i<\omega}\rangle$ and an automorphism $\sigma: K\to K$. For each $m<\omega$ there is $n\geq m$ such that $\sigma(a_1),\dots,\sigma(a_m)\in k\langle a_1,\dots,a_n\rangle$. As before, we know that $\sigma(a_1),\dots,\sigma(a_m)\equiv a_1,\dots,a_m$. Also as before, we have $(\sigma(a_1),\dots,\sigma(a_m))\in\dcl(a_1,\dots,a_n)$ so that there is a definable map $F:R^n\to R^m$ with $F(a_1,\dots,a_n)=(\sigma(a_1),\dots,\sigma(a_m))$. To show that $\sigma$ is the identity it suffices to show that $\sigma$ fixes $(a_1,\dots,a_m)$ for arbitrarily large $m$. Then we wish to choose $(a_i)_{i<\omega}$ to satisfy the following: for arbitrarily large $m$, for all $n\geq m$, given a definable map $F:R^n\to R^m$, either $F(a_1,\dots,a_n)=(a_1,\dots,a_m)$ or else $F(a_1,\dots,a_n)\not\equiv (a_1,\dots,a_m)$. This suggests the appropriate generalization of Lemma $\ref{MainLemma}$:

\begin{lemma}\label{ProjLemma}
    For any definable map $F:R^n\to R^m$ and any end cell $C\subseteq R^n$, there exists an end cell $C'\subseteq C$ such that either $F\upharpoonright C'=\pi_{\leq m}$, or else $\pi_{\leq m}(C')\cap F(C')=\emptyset$.
\end{lemma}

We will again postpone the proof to Section 3 and use Lemma \ref{ProjLemma} to prove the second main result; the particular construction implies Theorem \ref{FiniteThm} as well.



\begin{theorem}\label{InfiniteThm}
    There exists a non-Archimedean rigid real closed field of transcendence degree $\omega$ over $k$.
\end{theorem}

\begin{proof}
    Over all pairs $2\leq m\leq n$, enumerate all $k$-definable maps from $R^n\to R^m$, as $F_0,F_1,\dots$. We construct a partial type $\tau((x_i)_{i<\omega})$ in countably many variables as an intersection of a countable decreasing sequence of definable subsets $C_0\supseteq C_1\supseteq\dots$ of $R^\omega$. More specifically, each $C_i$, say defined by a formula $\varphi_i(x_1,\dots,x_{N_i})$, will be an end-cell when viewed as a definable subset of $R^{N_i}$. We will have $\varphi_{i+1}\to\varphi_{i}$ for all $i$ and we will take $\tau$ to be $\{\varphi_i(x_1,\dots,x_{N_i})\}_{i<\omega}$. Take $C_0$ to be the definable set $\{x_1>0\}$. Supposing $C_i$ has been defined and is an end cell in $R^{N_i
    }$, consider $F_i:R^n\to R^m$, $n\geq m\geq 2$. If $N_i<n$ then let $D$ be the $n$-dimensional end-cell $C_i\times R^{n-N_i}$. Now let $C'_{i+1}$ be the cell supplied by Lemma \ref{ProjLemma} applied to $D$ and $F_i$. Take $C_{i+1}$ to be $C'_{i+1}\cap \{x_1>i+1\}$. Otherwise if $N_i\geq n$ let $D$ be $\pi_{\leq n}(C_i)$ which is an end-cell of dimension $n$. Take $C'_{i+1}$ to be the end-cell supplied by Lemma \ref{ProjLemma} applied to $D$ and $F_i$. Then $C_i\cap \pi^{-1}_{\leq i}(C'_{i+1})$ is an end-cell contained in $C_i$. Take $C_{i+1}$ to be $C_i\cap \pi^{-1}_{\leq i}(C'_{i+1})\cap\{x_1>i+1\}$. (It is worth noting that the splitting into two cases is really non-essential here and is only mentioned for the sake of visualization; the two cases are identical if we think only of the defining formulas involved and the subsets of $R^\omega$ which they define). We may view $C_{i+1}$ as an end-cell in $R^{N_{i+1}}$ where $N_{i+1}:=\max\{n,N_i\}$. Now we have either that $F(x_1,\dots,x_n)=(x_1,\dots,x_m)$ for any $(x_1,\dots, x_{N_{i+1}})\in C_{i+1}$, or else for all $(x_1,\dots, x_{N_i+1})\in C_{i+1}$, there is no point $(y_1,\dots,y_{N_{i+1}})\in C_{i+1}$ with $(y_1,\dots,y_m)=F(x_1,\dots,x_n)$
    
    Now let $\tau((x_i)_{i<\omega})$ be the partial type in countable many variables given by $\bigcap_{i<\omega} C_i$. Then $\tau((x_i)_{i<\omega})\vdash x_1>\alpha$ for all $\alpha\in k$ and for any definable map $F:R^n\to R^m$, either $\tau((x_i)_{i<\omega})\vdash F(x_1,\dots,x_n)=(x_1,\dots,x_m)$ or else there is $\varphi(x_1,\dots,x_m)\in \tau((x_i)_{i<\omega})$ such that $\tau((x_i)_{i<\omega})\vdash \lnot\varphi(F(x_1,\dots,x_n))$. A little more precisely, there is a formula $\psi(x_1,\dots,x_n)\in\tau((x_i)_{i<\omega})$ such that $\psi(x_1,\dots,x_n)\vdash \lnot \exists y_{m+1}\dots y_n\psi(F(x_1),\dots,F(x_m),y_{m+1},\dots,y_n)$. In fact, if $F=F_i$, then $\psi$ may be taken as the formula defining $C_{i+1}'$ in the above construction.

    Since each $C_i$ is an end-cell, and in particular has full dimension as a subset of $R^{N_i}$, we may argue as in the finite case to take $(a_i)_{i<\omega}$ a realization in $R^\omega$ of $\tau(x_i)_{i<\omega}$ such that $(a_i)_{i<\omega}$ are algebraically independent. Now consider the real closure $K:=k\langle (a_i)_{i<\omega}\rangle$. This field is non-Archimedean as $a_1>\alpha$ for all $\alpha\in k$. Let $\sigma:K\to K$ be some automorphism of $K$. To show that $\sigma$ is the identity map, it suffices to show that $\sigma$ fixes the tuple $(a_1,\dots,a_m)$ for arbitrarily large $m$. In fact, we will show this is true for all $m\geq 2$.

    The tuple $(\sigma(a_1),\dots,\sigma(a_m))\in k\langle a_1,\dots,a_n\rangle$ for some sufficiently large $n\geq m$. Then there must be some definable map $F:R^n\to R^m$ with $F(a_1,\dots,a_n)=(\sigma(a_1),\dots,\sigma(a_m))$. Since $\sigma(a_1),\dots,\sigma(a_m)$ also satisfies the restriction of $\tau$ to the first $m$ variables, the construction of $\tau$ requires that $\tau((x_i)_{i<\omega})\vdash F(x_1,\dots,x_n)=(x_1,\dots,x_m)$. Hence $(\sigma(a_1),\dots,\sigma(a_m))=(a_1,\dots,a_m)$, which is what we needed to show.
\end{proof}


\end{document}
